% Recommended manuscript changes (static guidance; numbers in the tables).
\paragraph{Main paper.} Replace the single-target showcase with two pieces of
evidence that referees ask for: (i) a \emph{scaling} result and (ii) a
\emph{realistic} benchmark.
\begin{itemize}
\item Promote the \textbf{MoG mode-count scaling} figure (Fig.~\ref{fig:mog_scaling})
and the \textbf{ManyWell dimension-scaling} figure (Fig.~\ref{fig:mw_scaling}) to
the main paper: they show how coverage / EMC degrade with $K$ and $d$ for the
local baselines while LSB-MC stays close to the reference.
\item Add the \textbf{Bayesian GMM label-switching} experiment
(Fig.~\ref{fig:bg_modeweights}, Table~\ref{tab:bayes_gmm}) as the headline
``real model'' result: the jump bank is derived from a genuine symmetry of the
posterior rather than hand-tuned, which strengthens the methodological story.
\item Include the strongest baselines (PT and HMC), not just ULA/MALA/FLMC. The
honest comparison against PT is what makes the multimodal claim credible.
\item Show one \textbf{compute-fairness} panel (quality vs.\ wall-clock /
gradient evaluations) so the speed-up is not attributed to extra compute.
\end{itemize}

\paragraph{Appendix.} Per-dimension and per-$K$ tables, acceptance/swap
diagnostics, the exact configurations, and the reference-construction details for
the Bayesian posterior.

\paragraph{What the old section should become.} The original double-well /
single-MoG / single-ManyWell demonstrations are good \emph{illustrations} and
should move to the appendix or a ``setup'' figure; the quantitative claims should
rest on the scaling studies and the Bayesian benchmark.
